\begingroup\scriptsize\setlength{\tabcolsep}{3pt}
\begin{longtable}{@{}p{0.52in}p{0.85in}p{2.5in}p{0.22in}p{0.42in}p{1.95in}@{}}
\toprule
\textbf{ID} & \textbf{Target} & \textbf{Item} & \textbf{Tr} & \textbf{Prio} & \textbf{Sources} \\
\midrule
\endfirsthead
\toprule
\textbf{ID} & \textbf{Target} & \textbf{Item} & \textbf{Tr} & \textbf{Prio} & \textbf{Sources} \\
\midrule
\endhead
\midrule\multicolumn{6}{r}{\emph{continued}}\\\endfoot
\bottomrule\endlastfoot
N-a-001 & P3 3.5 countable-models & \(\omega\)-categoricity: Ryll-Nardzewski/Engeler/Svenonius characterization via types, DLO and random-graph examples & E & must & Ryll-Nardzewski 1959; Engeler 1959; Svenonius 1959; Hodges 1993 \\
N-a-002 & P3 3.5 countable-models & Fraisse constructions: hereditary/joint-embedding/amalgamation properties, existence and uniqueness of the Fraisse limit, random graph and extension axioms & E & must & Fraisse 1954; Rado 1964; Hodges 1993 \\
N-a-003 & P3 3.5 countable-models & prime, atomic, countably saturated, and homogeneous models: existence and uniqueness theorems over countable complete theories & E & must & Vaught 1961; Chang-Keisler 1990; Marker 2002 \\
N-a-004 & P3 3.5 countable-models & Vaught's never-two theorem (T); Vaught conjecture status report (P) & T & should & Vaught 1961; Morley 1970 \\
N-a-005 & P3 3.5 countable-models & two-cardinal theorems: Vaught's two-cardinal theorem statement and method sketch & S & could & Vaught 1965; Chang-Keisler 1990 \\
N-a-006 & P3 3.6 stability-classification & Morley rank and totally transcendental theories; outline of the proof of Morley's categoricity theorem; Baldwin-Lachlan refinement statement & S & should & Morley 1965; Baldwin-Lachlan 1971; Marker 2002 \\
N-a-007 & P3 3.6 stability-classification & stability spectrum and forking: stable/superstable/\(\omega\)-stable, nonforking extensions, symmetry and transitivity survey & S & should & Shelah 1978; Pillay 1996; Tent-Ziegler 2012 \\
N-a-008 & P3 3.6 stability-classification & indiscernible sequences and Ehrenfeucht-Mostowski models: existence via Ramsey's theorem, applications to non-isomorphic models & T & should & Ehrenfeucht-Mostowski 1956; Ramsey 1930; Hodges 1993 \\
N-a-009 & P3 3.7 tame-topology & o-minimality: definition, monotonicity and cell-decomposition statements, RCF as an o-minimal structure & S & should & Pillay-Steinhorn 1986; Knight-Pillay-Steinhorn 1986; van den Dries 1998 \\
N-a-010 & P3 3.7 tame-topology & NIP and simple theories: definitions, examples, map of the neostability landscape & S & could & Shelah 1971; Simon 2015; Tent-Ziegler 2012 \\
N-a-011 & P3 3.8 preservation-model-completeness & preservation theorems: Los-Tarski (universal), Lyndon (positivity), Chang-Los-Suszko (unions of chains) & T & should & Los 1955; Lyndon 1959; Chang-Keisler 1990; Hodges 1993 \\
N-a-012 & P3 3.8 preservation-model-completeness & model completeness and model companions: Robinson's test, ACF and RCF worked examples, existentially closed structures & E & must & Robinson 1956; Chang-Keisler 1990; Marker 2002 \\
N-a-013 & P3 3.4 interpolation-definability & Robinson joint consistency theorem and Beth definability with model-theoretic proofs; equivalence with Craig interpolation & T & should & Robinson 1956; Beth 1953; Hodges 1993 \\
N-a-014 & P3 3.2 types-saturation & ultrafilter existence from BPI, BPI versus AC status, Keisler-Shelah ultrapower isomorphism theorem statement & S & could & Keisler 1961; Shelah 1971; Bell-Slomson 1969 \\
N-a-015 & P7 7.4 lambda-calculi & untyped lambda calculus and combinatory logic as computation models: Church-Rosser theorem statement, lambda-definability equals partial recursiveness & T & must & Church 1936; Church-Rosser 1936; Kleene 1936; Curry-Feys 1958; Barendregt 1984 \\
N-a-016 & P4 4.4 degrees-and-priority & Turing degrees: partial order, join, jump operator, Kleene-Post construction of incomparable degrees below \(\mathbf 0'\) & E & must & Kleene-Post 1954; Rogers 1967; Soare 2016 \\
N-a-017 & P4 4.4 degrees-and-priority & Post's problem and the Friedberg-Muchnik theorem: finite-injury priority construction in full & T & must & Post 1944; Friedberg 1957; Muchnik 1956; Soare 2016 \\
N-a-018 & P4 4.4 degrees-and-priority & simple and creative sets, many-one and one-one reducibility, Myhill isomorphism theorem & T & should & Post 1944; Myhill 1955; Rogers 1967 \\
N-a-019 & P4 4.5 frontiers-of-computability & hyperarithmetical and analytical hierarchies: \(\Delta^1_1=\mathrm{HYP}\) statement, ordinal notations \(\mathcal O\) sketch & S & should & Kleene 1955; Rogers 1967; Sacks 1990 \\
N-a-020 & P4 4.5 frontiers-of-computability & Kolmogorov complexity: invariance theorem, incompressibility, Chaitin-style incompleteness & T & should & Kolmogorov 1965; Chaitin 1974; Li-Vitanyi 2008; Downey-Hirschfeldt 2010 \\
N-a-021 & P4 4.5 frontiers-of-computability & Hilbert's tenth problem: MRDP theorem statement, Diophantine sets equal c.e. sets, proof survey with Pell-equation core & S & must & Davis-Putnam-Robinson 1961; Matiyasevich 1970; Davis 1973 \\
N-a-022 & P4 4.5 frontiers-of-computability & undecidable word problems: Post-Markov theorem for semigroups (T), Novikov-Boone theorem for groups (S) & T & could & Post 1947; Markov 1947; Novikov 1955; Boone 1958 \\
N-a-023 & P4 4.5 frontiers-of-computability & busy beaver function: definition and Rado's noncomputability theorem & T & could & Rado 1962 \\
N-a-024 & P4 4.5 frontiers-of-computability & computable structure theory pointer: computable presentations, indices of structures & P & could & Ash-Knight 2000 \\
N-a-025 & P5 5.6 provability-logic & GL: axioms, Kripke semantics for transitive conversely-well-founded frames, completeness and FMP, fixed-point theorem, Solovay arithmetical completeness (S) & E & must & Solovay 1976; Sambin 1976; Segerberg 1971; Boolos 1993 \\
N-a-026 & P5 5.6 provability-logic & Rosser provability and witness comparison; Guaspari-Solovay \(R\) pointer & P & could & Guaspari-Solovay 1979; Lindstrom 2003 \\
N-a-027 & P5 5.5 independence-and-strength & Paris-Harrington theorem: statement, unprovability in PA via indicator/model-theoretic argument (S), provability in full-induction ACA and nonprovability in ACA_0 & S & must & Paris-Harrington 1977; Kanamori-McAloon 1987; Kaye 1991 \\
N-a-028 & P5 5.5 independence-and-strength & Goodstein sequences and the Kirby-Paris hydra: termination via \(\varepsilon_0\) (T), unprovability in PA statement & T & must & Goodstein 1944; Kirby-Paris 1982 \\
N-a-029 & P5 5.5 independence-and-strength & interpretability and relative interpretation: Tarski-Mostowski-Robinson method, essential undecidability transfer, interpretability logic pointer & S & should & Tarski-Mostowski-Robinson 1953; Visser 1997; Lindstrom 2003 \\
N-a-030 & P5 5.5 independence-and-strength & subsystems of second-order arithmetic and reverse mathematics: \(\mathsf{RCA}_0\), \(\mathsf{WKL}_0\), \(\mathsf{ACA}_0\), \(\mathsf{ATR}_0\), \(\Pi^1_1\text{-}\mathsf{CA}_0\); Big Five phenomenon with sample equivalences & S & must & Friedman 1975; Simpson 2009 \\
N-a-031 & P5 5.5 independence-and-strength & intensionality of consistency statements: Feferman's analysis, dependence on the provability predicate & S & should & Feferman 1960 \\
N-a-032 & P5 5.4 models-of-pa & nonstandard models of PA: existence and order type \(\mathbb N+\mathbb Z\cdot\mathbb Q\), overspill and its uses, Tennenbaum's theorem & E & must & Skolem 1934; Tennenbaum 1959; Kaye 1991 \\
N-a-033 & P5 5.7 decidable-arithmetic & Presburger arithmetic: quantifier elimination with congruence predicates, decidability, Fischer-Rabin lower bound statement & T & must & Presburger 1929; Fischer-Rabin 1974; Enderton 2001 \\
N-a-034 & P5 5.7 decidable-arithmetic & Skolem arithmetic of multiplication: decidability via Feferman-Vaught statement & S & could & Skolem 1930; Feferman-Vaught 1959 \\
N-a-035 & P6 6.5 cardinal-arithmetic & cofinality, regular and singular cardinals, Konig's theorem, Hausdorff recursion, Easton's theorem statement & E & must & Konig 1905; Easton 1970; Jech 2003; Kunen 2011 \\
N-a-036 & P6 6.5 cardinal-arithmetic & singular cardinal behavior: Silver's theorem statement, PCF pointer & S & could & Silver 1975; Shelah 1994 \\
N-a-037 & P6 6.4 infinitary-combinatorics & closed unbounded and stationary sets, Fodor's pressing-down lemma, diagonal intersection & E & must & Fodor 1956; Jech 2003; Kunen 2011 \\
N-a-038 & P6 6.4 infinitary-combinatorics & Delta-system lemma with proof and the ccc verification pattern for Cohen forcing & T & must & Shanin 1946; Kunen 2011 \\
N-a-039 & P6 6.4 infinitary-combinatorics & infinite Ramsey theorem (T); Erdos-Rado theorem and partition calculus notation (S) & T & must & Ramsey 1930; Erdos-Rado 1956; Jech 2003 \\
N-a-040 & P6 6.4 infinitary-combinatorics & trees: Konig infinity lemma, Aronszajn tree construction, Suslin hypothesis and its independence status survey & S & should & Konig 1927; Kurepa 1935; Todorcevic 1984; Kunen 2011 \\
N-a-041 & P6 6.4 infinitary-combinatorics & combinatorial principles: diamond and square statements with one application each; Martin's axiom with ccc applications & S & should & Jensen 1972; Martin-Solovay 1970; Kunen 2011 \\
N-a-042 & P6 6.3 constructibility-forcing & Cohen forcing of \(\neg\mathrm{CH}\): full ccc cardinal-preservation verification, nice names counting & T & must & Cohen 1963; Kunen 2011 \\
N-a-043 & P6 6.3 constructibility-forcing & independence of AC: symmetric extensions method survey; permutation models of ZFA & S & should & Fraenkel 1922; Cohen 1963; Jech 1973; Halbeisen 2017 \\
N-a-044 & P6 6.6 large-cardinals & large cardinals: inaccessible, Mahlo, weakly compact, measurable via ultrapowers and elementary embeddings, Scott's theorem \(V\neq L\), consistency-strength hierarchy chart to supercompact & S & must & Ulam 1930; Scott 1961; Keisler 1962; Kanamori 2003 \\
N-a-045 & P6 6.6 large-cardinals & Levy hierarchy, absoluteness for transitive models, Shoenfield absoluteness statement with applications & T & should & Levy 1965; Shoenfield 1961; Jech 2003 \\
N-a-046 & P6 6.7 descriptive-set-theory & descriptive set theory: Polish spaces, Borel hierarchy, analytic sets, Suslin's theorem, perfect-set property, regularity properties & E & must & Suslin 1917; Kechris 1995; Moschovakis 2009 \\
N-a-047 & P6 6.7 descriptive-set-theory & effective descriptive set theory bridge to computability: lightface pointclasses pointer & P & could & Moschovakis 2009; Sacks 1990 \\
N-a-048 & P6 6.8 determinacy & infinite games: Gale-Stewart theorem (T), Borel determinacy statement, AD versus AC, projective determinacy and Woodin cardinals pointer & S & should & Gale-Stewart 1953; Mycielski-Steinhaus 1962; Martin 1975; Martin-Steel 1989 \\
N-a-049 & P7 7.5 arithmetic-proof-theory & Gentzen's consistency proof for PA: reduction procedure on derivations, transfinite induction up to \(\varepsilon_0\), ordinal assignment in full for the propositional core & T & must & Gentzen 1936; Gentzen 1938; Takeuti 1987; Schwichtenberg 1977 \\
N-a-050 & P7 7.5 arithmetic-proof-theory & ordinal notation systems: Cantor normal form arithmetic, Veblen functions, \(\Gamma_0\) and Feferman-Schutte statement & S & should & Veblen 1908; Feferman 1964; Schutte 1977; Pohlers 2009 \\
N-a-051 & P7 7.5 arithmetic-proof-theory & PRA and finitism: Skolem's arithmetic, Tait's identification of finitism with PRA & S & should & Skolem 1923; Tait 1981 \\
N-a-052 & P7 7.5 arithmetic-proof-theory & omega-logic: the \(\omega\)-rule, cut elimination for PA with the \(\omega\)-rule, bound \(\varepsilon_0\) & S & could & Schutte 1951; Schwichtenberg 1977; Pohlers 2009 \\
N-a-053 & P7 7.6 functional-interpretations & Godel's Dialectica interpretation: System T, interpretation of HA, consistency corollary & S & should & Godel 1958; Troelstra 1973; Avigad-Feferman 1998 \\
N-a-054 & P7 7.6 functional-interpretations & Kleene number realizability: soundness for HA, ECT and Markov's principle independence applications & T & must & Kleene 1945; Troelstra 1998; van Oosten 2008 \\
N-a-055 & P7 7.6 functional-interpretations & disjunction and numerical existence properties for IPC and HA with proofs & T & must & Godel 1932; Gentzen 1935; Kleene 1945; Troelstra-van Dalen 1988 \\
N-a-056 & P7 7.7 structural-extensions & epsilon calculus: epsilon terms, first and second epsilon theorems statements, Herbrand's theorem in general form & S & should & Hilbert-Bernays 1939; Herbrand 1930; Buss 1998 \\
N-a-057 & P7 7.7 structural-extensions & generalized sequent frameworks survey: hypersequents, labelled calculi, display calculi, deep inference, with one worked modal example each & S & should & Avron 1996; Negri 2005; Belnap 1982; Guglielmi 2007; Negri-von Plato 2001 \\
N-a-058 & P7 7.7 structural-extensions & proof mining pointer: extraction of bounds from classical proofs & P & could & Kohlenbach 2008 \\
N-a-059 & P3 3.9 abstract-model-theory & infinitary logic \(L_{\omega_1\omega}\): syntax, Scott isomorphism theorem, Scott sentences, Karp completeness statement, compactness failure & S & should & Karp 1964; Scott 1965; Keisler 1971; Marker 2016 \\
N-a-060 & P3 3.9 abstract-model-theory & generalized quantifiers: Mostowski and Lindstrom quantifiers, completeness for \(L(Q_1)\) statement & S & could & Mostowski 1957; Lindstrom 1966; Keisler 1970 \\
N-a-061 & P3 3.9 abstract-model-theory & abstract logics and the full proof of Lindstrom's first theorem & T & should & Lindstrom 1969; Barwise 1974; Ebbinghaus-Flum-Thomas 1994 \\
N-a-062 & P3 3.9 abstract-model-theory & independence-friendly and dependence logic pointer & P & could & Hintikka-Sandu 1989; Vaananen 2007 \\
N-a-063 & P4 4.1 machines-and-models & \(\mu\)-recursion and partial recursive functions: normal form and s-m-n recap, equivalence theorem Turing-computable \(\equiv\) \(\mu\)-recursive & E & must & Kleene 1936; Turing 1936; Kleene 1952; Odifreddi 1989 \\
N-a-064 & P4 4.6 complexity & complexity theory as a subject: class ladder P, NP, coNP, PH, PSPACE, EXPTIME, time and space hierarchy theorems, logic-relevant completeness results table & T & must & Hartmanis-Stearns 1965; Cook 1971; Karp 1972; Stockmeyer 1976; Papadimitriou 1994; Arora-Barak 2009 \\
N-a-065 & P3 3.1 elementary-maps-ef-games & Ehrenfeucht-Fraisse games: game characterization of elementary equivalence, quantifier rank, inexpressibility applications & E & must & Fraisse 1954; Ehrenfeucht 1961; Hodges 1993; Ebbinghaus-Flum-Thomas 1994 \\
N-a-066 & P6 6.1 ordinal-arithmetic & ordinal arithmetic developed as theory: sums, products, exponentiation, Cantor normal form, transfinite induction and recursion theorems & E & must & Cantor 1897; Sierpinski 1965; Jech 2003; Kunen 2011 \\
N-a-067 & P6 6.3 constructibility-forcing & constructible hierarchy \(L\): condensation lemma statement, \(\mathrm{GCH}\) in \(L\), reflection principle, fine-structure pointer & S & should & Godel 1938; Levy 1960; Devlin 1984; Jech 2003 \\
N-a-068 & P6 6.3 constructibility-forcing & Boolean-valued models at statement level; product and iterated forcing pointers & S & should & Scott-Solovay 1967; Bell 2005; Kunen 2011 \\
N-a-069 & P7 7.2 cut-elimination & midsequent theorem via cut elimination and Herbrand's theorem as corollary & S & should & Gentzen 1935; Herbrand 1930; Takeuti 1987; Buss 1998 \\
N-a-070 & P7 7.2 cut-elimination & Craig interpolation proved by Maehara's method with Beth definability corollary & E & must & Craig 1957; Maehara 1960; Takeuti 1987 \\
N-a-071 & P2 2.3 second-higher-order & second-order and HOL proof calculi displayed: comprehension and choice schemas as a calculus, Henkin completeness route & E & must & Church 1940; Henkin 1950; Andrews 2002; Shapiro 1991 \\
N-b-001 & P8 8.2 modal & canonical models and completeness for K, T, D, B, K4, S4, S5: Lindenbaum construction for normal modal logics, canonicity discussion & E & must & Lemmon-Scott 1977; Chellas 1980; Blackburn-de Rijke-Venema 2001 \\
N-b-002 & P8 8.2 modal & frame correspondence: definability of the T/4/5/B/D axioms, Sahlqvist theorem statement with algorithmic examples (proof S) & T & must & Sahlqvist 1975; van Benthem 1983; Blackburn-de Rijke-Venema 2001 \\
N-b-003 & P8 8.2 modal & bisimulation: definition, invariance, van Benthem characterization theorem & T & must & van Benthem 1976; Blackburn-de Rijke-Venema 2001 \\
N-b-004 & P8 8.2 modal & filtration, finite model property, and decidability for K, T, S4, S5 & T & must & Segerberg 1971; Chagrov-Zakharyaschev 1997; Blackburn-de Rijke-Venema 2001 \\
N-b-005 & P8 8.2 modal & Kripke frame incompleteness phenomena: general frames, Thomason and Fine counterexamples survey & S & should & Thomason 1974; Fine 1974; Chagrov-Zakharyaschev 1997 \\
N-b-006 & P8 8.2 modal & non-normal modal logics and neighborhood semantics: systems E, M, C, N, regular logics, completeness & E & should & Scott 1970; Montague 1970; Chellas 1980; Pacuit 2017 \\
N-b-007 & P8 8.2 modal & per-system entry cards for the modal cube: K, D, T, B, K4, S4, S4.2, S5, GL, Grz -- axioms, frame class, FMP, decidability, characteristic theorems & E & must & Chellas 1980; Hughes-Cresswell 1996; Chagrov-Zakharyaschev 1997 \\
N-b-008 & P8 8.2 modal & two-dimensional semantics: actuality operator, fixedly, Kaplan character/content distinction & S & should & Crossley-Humberstone 1977; Davies-Humberstone 1980; Kaplan 1989 \\
N-b-009 & P8 8.2 modal & quantified modal logic: constant versus varying domains, Barcan and converse Barcan formulas, rigid designators, existence predicate, actualism/possibilism & E & must & Barcan 1946; Kripke 1963; Fitting-Mendelsohn 1998; Garson 2001; Williamson 2013 \\
N-b-010 & P8 8.4 conditionals-deontic-free & free logic: positive/negative/neutral systems, inner and outer domains, definite descriptions (Russell versus free-logic treatments) & E & should & Russell 1905; Lambert 1967; Bencivenga 1986; Lehmann 2002 \\
N-b-011 & P8 8.3 temporal-dynamic-epistemic & multi-agent epistemic logic: S5\(_n\)/KD45\(_n\), common knowledge as fixed point, muddy children analysis & E & must & Hintikka 1962; Aumann 1976; Fagin-Halpern-Moses-Vardi 1995 \\
N-b-012 & P8 8.3 temporal-dynamic-epistemic & dynamic epistemic logic: public announcement logic axioms, reduction laws, unsuccessful updates & S & should & Plaza 1989; Baltag-Moss-Solecki 1998; van Ditmarsch-van der Hoek-Kooi 2008 \\
N-b-013 & P8 8.4 conditionals-deontic-free & deontic logic: SDL as KD, Ross and Good Samaritan and Chisholm contrary-to-duty paradoxes, dyadic conditional obligation & E & must & von Wright 1951; Chisholm 1963; Hansson 1969; McNamara 2006 \\
N-b-014 & P8 8.3 temporal-dynamic-epistemic & philosophical tense logic: Prior's systems, since/until and Kamp's expressive completeness statement, Ockhamist versus Peircean branching time & S & should & Prior 1967; Kamp 1968; Burgess 1984 \\
N-b-015 & P8 8.3 temporal-dynamic-epistemic & hybrid logic: nominals, satisfaction operators, axiomatization survey & S & could & Prior 1967; Blackburn 2000; Areces-ten Cate 2007 \\
N-b-016 & P8 8.4 conditionals-deontic-free & counterfactuals: Stalnaker selection functions versus Lewis sphere systems, conditional excluded middle debate, limit assumption, completeness statements for C1 and VC & E & must & Stalnaker 1968; Lewis 1973; Nute 1980 \\
N-b-017 & P8 8.4 conditionals-deontic-free & indicative conditionals: paradoxes of material implication, Adams's thesis, Lewis triviality results, suppositional theories & S & must & Adams 1975; Lewis 1976; Jackson 1987; Edgington 1995; Bennett 2003 \\
N-b-018 & P8 8.4 conditionals-deontic-free & conditional logics as a family: CK and neighbors, selection-function and premise semantics, correspondence table & T & should & Chellas 1975; Nute 1980 \\
N-b-019 & P11 11.2 truth-theories & Tarskian truth as a philosophical unit: convention T, object/metalanguage discipline, hierarchy of languages & S & should & Tarski 1935; Tarski 1944 \\
N-b-020 & P11 11.2 truth-theories & Kripke's fixed-point theory of truth: partial Kleene valuation schemes, least and maximal fixed points, groundedness, ghost of the Tarski hierarchy & T & must & Kripke 1975; Martin 1984; Halbach 2011 \\
N-b-021 & P11 11.2 truth-theories & revision theory of truth survey & S & could & Herzberger 1982; Gupta-Belnap 1993 \\
N-b-022 & P11 11.2 truth-theories & axiomatic theories of truth: TB/CT conservativity, KF, FS, comparative survey & S & should & Friedman-Sheard 1987; Feferman 1991; Halbach 2011; Horsten 2011 \\
N-b-023 & P11 11.1 paradoxes & paradox catalog with formal diagnoses: strengthened liar, Curry, Yablo, Grelling, Berry & E & must & Curry 1942; Yablo 1993; Priest 2006; Beall-Glanzberg-Ripley 2018 \\
N-b-024 & P11 11.4 vagueness & vagueness: sorites structure, supervaluationism with global versus local validity, degree-theoretic approaches, epistemicism, contextualism pointer & S & must & Fine 1975; Williamson 1994; Keefe 2000; Smith 2008 \\
N-b-025 & P8 8.5 many-valued-fuzzy & first-degree entailment: Belnap-Dunn four-valued semantics, bilattice reading, relation to K3 and LP & E & must & Dunn 1976; Belnap 1977; Omori-Wansing 2017 \\
N-b-026 & P8 8.7 paraconsistent-nonmonotonic & dialetheism and the choice of paraconsistent consequence: motivation from paradox, inclosure schema, replies & S & should & Priest 1979; Priest 2006 \\
N-b-027 & P8 8.6 substructural-linear & relevance logic landscape: E and R axiomatic systems, variable-sharing proof, ternary-frame conditions per axiom, admissibility of gamma & S & should & Anderson-Belnap 1975; Meyer-Dunn 1969; Routley-Meyer 1973; Dunn-Restall 2002; Mares 2004 \\
N-b-028 & P8 8.5 many-valued-fuzzy & Lukasiewicz logics: L3, finite Ln, infinite-valued L, axiomatization, McNaughton theorem statement, real-valued semantics & E & must & Lukasiewicz 1920; Lukasiewicz-Tarski 1930; McNaughton 1951; Cignoli-D'Ottaviano-Mundici 2000 \\
N-b-029 & P8 8.5 many-valued-fuzzy & t-norm fuzzy logics: continuous t-norms, BL and MTL, Godel and product logics, standard completeness survey & S & must & Hajek 1998; Esteva-Godo 2001; Cintula-Hajek-Noguera 2011 \\
N-b-030 & P8 8.1 intuitionistic-intermediate & intuitionism beyond the calculus: BHK reading, Brouwerian counterexamples, choice sequences, bar induction, continuity principles, Markov's principle status & S & should & Brouwer 1908; Heyting 1956; Dummett 1977; Troelstra-van Dalen 1988 \\
N-b-031 & P8 8.1 intuitionistic-intermediate & intermediate logics: Godel-Dummett LC, KC, the lattice of superintuitionistic logics survey & S & should & Godel 1932; Dummett 1959; Jankov 1968; Chagrov-Zakharyaschev 1997 \\
N-b-032 & P8 8.7 paraconsistent-nonmonotonic & nonmonotonic consequence: preferential models, KLM system P, rational monotony; AGM belief revision postulates, Levi and Harper identities & E & must & Alchourron-Gardenfors-Makinson 1985; Kraus-Lehmann-Magidor 1990; Lehmann-Magidor 1992; Makinson 2005 \\
N-b-033 & P8 8.7 paraconsistent-nonmonotonic & default logic and circumscription entry cards & S & should & Reiter 1980; McCarthy 1980 \\
N-b-034 & P11 11.3 consequence-and-constants & philosophy of logical consequence: Tarski's analysis, model-theoretic versus proof-theoretic accounts, demarcation of logical constants (invariance), harmony & S & should & Tarski 1936; Prawitz 1974; Etchemendy 1990; Sher 1991; Dummett 1991 \\
N-b-035 & P11 11.5 ontological-commitment & plural quantification and the status of second-order logic: Boolos's plural reading, ontological commitment, Quine's 'set theory in disguise' charge & S & should & Quine 1970; Boolos 1984; Boolos 1985; Shapiro 1991 \\
N-c-001 & P9 9.1 finite-model-theory & finite model theory foundations: failure of compactness and completeness in the finite, Craig/Beth failure, EF-game method on finite structures, Trakhtenbrot recap & E & must & Ebbinghaus-Flum 1995; Libkin 2004 \\
N-c-002 & P9 9.1 finite-model-theory & zero-one laws for first-order logic: Glebskii and Fagin theorems, extension axioms, the random-graph link & T & must & Glebskii et al. 1969; Fagin 1976; Libkin 2004 \\
N-c-003 & P9 9.1 finite-model-theory & locality: Hanf and Gaifman theorems with inexpressibility applications (connectivity, acyclicity) & T & should & Hanf 1965; Gaifman 1982; Libkin 2004 \\
N-c-004 & P9 9.1 finite-model-theory & descriptive complexity: Fagin's theorem (existential SO captures NP) with proof sketch, Immerman-Vardi (LFP captures P on ordered structures) statement & T & must & Fagin 1974; Vardi 1982; Immerman 1986; Immerman 1999 \\
N-c-005 & P9 9.2 decidable-fragments & the classical decision problem: prefix-class classification, Bernays-Schonfinkel-Ramsey, Ackermann, Godel-Kalmar-Schutte classes, reduction classes & S & should & Lewis 1979; Borger-Gradel-Gurevich 1997 \\
N-c-006 & P9 9.2 decidable-fragments & two-variable fragment: finite model property and NEXPTIME decidability statement & S & should & Mortimer 1975; Gradel-Kolaitis-Vardi 1997 \\
N-c-007 & P9 9.2 decidable-fragments & guarded fragment: decidability and the modal connection & S & should & Andreka-van Benthem-Nemeti 1998; Gradel 1999 \\
N-c-008 & P9 9.2 decidable-fragments & monadic first-order logic: Lowenheim-Behmann decidability & T & should & Lowenheim 1915; Behmann 1922; Borger-Gradel-Gurevich 1997 \\
N-c-009 & P9 9.3 description-logics & ALC: syntax, model semantics, TBox/ABox, tableau decision procedure with blocking, ExpTime complexity statement & E & must & Schmidt-Schauss-Smolka 1991; Baader et al. 2003; Baader-Horrocks-Lutz-Sattler 2017 \\
N-c-010 & P9 9.3 description-logics & lightweight and expressive DLs: EL family with polynomial classification, DL-Lite and query answering, SHOIQ/OWL profiles map & S & should & Baader-Brandt-Lutz 2005; Calvanese et al. 2007; Baader-Horrocks-Lutz-Sattler 2017 \\
N-c-011 & P9 9.3 description-logics & ALC as notational variant of multimodal K (Schild correspondence) & T & should & Schild 1991 \\
N-c-012 & P9 9.4 temporal-logics-automata & LTL deepened: past operators, Kamp and Gabbay-Pnueli-Shelah-Stavi expressive completeness statements, safety/liveness, PSPACE satisfiability & T & must & Kamp 1968; Pnueli 1977; Gabbay-Pnueli-Shelah-Stavi 1980; Sistla-Clarke 1985 \\
N-c-013 & P9 9.4 temporal-logics-automata & omega-automata: Buchi/Rabin/parity acceptance, closure properties, complementation statement, Vardi-Wolper LTL-to-automata translation & S & must & Buchi 1962; McNaughton 1966; Vardi-Wolper 1986; Safra 1988; Thomas 1997 \\
N-c-014 & P9 9.4 temporal-logics-automata & MSO and automata: Buchi-Elgot-Trakhtenbrot theorem on finite words (T), S1S decidability, Rabin's S2S statement & T & should & Buchi 1960; Elgot 1961; Trakhtenbrot 1961; Rabin 1969; Thomas 1997 \\
N-c-015 & P9 9.4 temporal-logics-automata & tree automata pointer and applications to decidability transfer & P & could & Thatcher-Wright 1968; Comon et al. 2007 \\
N-c-016 & P9 9.5 model-checking & CTL and CTL*: syntax, semantics, fixpoint characterizations of EG/EU, labeling algorithm with complexity, LTL versus CTL expressiveness & E & must & Clarke-Emerson 1981; Queille-Sifakis 1982; Emerson-Halpern 1986; Clarke-Grumberg-Peled 1999; Baier-Katoen 2008 \\
N-c-017 & P9 9.5 model-checking & symbolic model checking: BDDs and Bryant canonicity, fixpoint computation, bounded model checking via SAT, CEGAR pointer & S & must & Bryant 1986; Burch et al. 1992; McMillan 1993; Biere-Cimatti-Clarke-Zhu 1999; Clarke et al. 2003 \\
N-c-018 & P9 9.5 model-checking & modal mu-calculus: syntax, Knaster-Tarski semantics, relation to parity games, alternation hierarchy statement & S & should & Kozen 1983; Emerson-Jutla 1991; Bradfield-Stirling 2007 \\
N-c-019 & P9 9.6 program-logics & Hoare logic: partial and total correctness, axioms and rules, soundness proof, relative completeness (Cook), weakest preconditions (Dijkstra) & E & must & Floyd 1967; Hoare 1969; Dijkstra 1975; Cook 1978; Apt 1981 \\
N-c-020 & P9 9.6 program-logics & separation logic: heap assertions, frame rule, symbolic execution, bi-abduction pointer & S & should & O'Hearn-Reynolds-Yang 2001; Reynolds 2002; Calcagno et al. 2011; O'Hearn 2019 \\
N-c-021 & P9 9.6 program-logics & propositional dynamic logic: syntax, regular programs, Segerberg axioms with Kozen-Parikh completeness statement, Fischer-Ladner closure, ExpTime decidability & E & should & Pratt 1976; Fischer-Ladner 1979; Kozen-Parikh 1981; Harel-Kozen-Tiuryn 2000 \\
N-c-022 & P9 9.7 logic-programming-databases & logic programming: definite programs, least Herbrand model, immediate-consequence fixpoint, SLD-resolution soundness and completeness, negation as failure with Clark completion, SLDNF & E & must & Kowalski 1974; van Emden-Kowalski 1976; Clark 1978; Apt-van Emden 1982; Lloyd 1987 \\
N-c-023 & P9 9.7 logic-programming-databases & answer set programming: stable-model semantics via Gelfond-Lifschitz reduct, complexity statement & S & should & Gelfond-Lifschitz 1988; Marek-Truszczynski 1999; Brewka-Eiter-Truszczynski 2011 \\
N-c-024 & P9 9.7 logic-programming-databases & Datalog and conjunctive queries: Chandra-Merlin homomorphism theorem, containment decidability, data/combined complexity & S & should & Chandra-Merlin 1977; Vardi 1982; Abiteboul-Hull-Vianu 1995 \\
N-c-025 & P9 9.7 logic-programming-databases & constraint satisfaction and logic pointer: CSP as homomorphism problem, dichotomy theorem statement & P & could & Feder-Vardi 1998; Bulatov 2017; Zhuk 2017 \\
N-c-026 & P9 9.8 complexity-of-logic & complexity of logical tasks: Cook-Levin recap, QBF PSPACE-completeness with proof, Ladner-style table for modal logics (K/S4 PSPACE, S5 NP), FO validity undecidable recap & T & must & Cook 1971; Stockmeyer-Meyer 1973; Ladner 1977; Papadimitriou 1994 \\
N-c-027 & P7 7.8 proof-complexity & proof complexity: Cook-Reckhow framework, resolution lower bound for the pigeonhole principle (Haken) proof sketch, Frege systems open-problems map & S & should & Cook-Reckhow 1979; Haken 1985; Beame-Pitassi 1998; Krajicek 2019 \\
N-c-028 & P9 9.8 complexity-of-logic & interpolation in verification pointer: McMillan-style interpolant-based model checking & P & could & McMillan 2003 \\
N-c-029 & P7 7.4 lambda-calculi & System F: syntax, impredicative encodings, Girard's strong normalization statement, Curry-Howard for second-order, parametricity pointer & S & must & Girard 1972; Reynolds 1974; Reynolds 1983; Girard-Lafont-Taylor 1989 \\
N-c-030 & P7 7.4 lambda-calculi & Martin-Lof type theory: judgments, Pi/Sigma/identity types, universes, propositions-as-types with proof relevance, normalization and decidable type checking statements & S & must & Martin-Lof 1984; Nordstrom-Petersson-Smith 1990; Hofmann 1997 \\
N-c-031 & P12 12.5 interactive-proof & lambda cube and pure type systems; calculus of (inductive) constructions as the Coq/Lean foundation & S & should & Coquand-Huet 1988; Barendregt 1992; Paulin-Mohring 2015 \\
N-c-032 & P12 12.1 sat & CDCL: clause learning with UIP, watched literals, restarts, VSIDS; GRASP/Chaff/MiniSat lineage; DRAT proofs and certified UNSAT & E & must & Marques-Silva-Sakallah 1999; Moskewicz et al. 2001; Een-Sorensson 2003; Marques-Silva-Lynce-Malik 2009; Wetzler-Heule-Hunt 2014 \\
N-c-033 & P12 12.1 sat & tractable cases and hardness structure: 2-SAT via implication graphs, Horn-SAT unit propagation, random k-SAT phase transition pointer & T & should & Aspvall-Plass-Tarjan 1979; Dowling-Gallier 1984; Mitchell-Selman-Levesque 1992 \\
N-c-034 & P12 12.2 smt & SMT internals: DPLL(T) architecture, congruence closure for EUF, simplex for linear arithmetic, Nelson-Oppen combination with proof for convex stably-infinite theories & E & must & Nelson-Oppen 1979; Dutertre-de Moura 2006; Nieuwenhuis-Oliveras-Tinelli 2006; de Moura-Bjorner 2008; Barrett-Sebastiani-Seshia-Tinelli 2009 \\
N-c-035 & P12 12.3 first-order-atp & unification theory deepened: Martelli-Montanari rule-based algorithm with termination proof, mgu theorem, occurs-check, equational and AC-unification survey, undecidability of higher-order unification & E & must & Robinson 1965; Huet 1973; Goldfarb 1981; Martelli-Montanari 1982; Baader-Snyder 2001 \\
N-c-036 & P12 12.3 first-order-atp & ordered resolution and the superposition calculus: reduction orders, redundancy and saturation, model-construction completeness proof sketch & S & must & Bachmair-Ganzinger 1994; Nieuwenhuis-Rubio 2001 \\
N-c-037 & P12 12.4 rewriting & term rewriting: abstract reduction systems, Newman's lemma, critical pairs and Knuth-Bendix completion, termination orders (LPO/KBO), word-problem applications & E & must & Newman 1942; Knuth-Bendix 1970; Huet 1980; Dershowitz 1982; Baader-Nipkow 1998 \\
N-c-038 & P12 12.5 interactive-proof & LCF architecture and kernels: de Bruijn criterion, proof terms versus proof scripts, trusted computing bases of Coq/Lean/Isabelle/HOL, reflection & S & should & de Bruijn 1970; Gordon-Milner-Wadsworth 1979; Pollack 1998; Harrison 2009 \\
N-c-039 & P12 12.5 interactive-proof & program extraction from constructive proofs: realizability-based extraction & S & could & Paulin-Mohring 1989; Letouzey 2003 \\
N-c-040 & P12 12.5 interactive-proof & ATP-ITP integration pointer: hammers and proof reconstruction & P & could & Paulson-Blanchette 2010; Blanchette et al. 2016 \\
N-d-001 & P10 10.1 universal-algebra & universal algebra toolkit: algebras, homomorphisms, congruences and quotients, subdirect representation, Birkhoff HSP variety theorem & T & must & Birkhoff 1935; Birkhoff 1944; Burris-Sankappanavar 1981 \\
N-d-002 & P10 10.1 universal-algebra & equational logic: term algebras, free algebras, Birkhoff completeness for equational consequence & T & must & Birkhoff 1935; Burris-Sankappanavar 1981 \\
N-d-003 & P10 10.2 boolean-stone & Boolean algebras: filters, ideals, ultrafilters, atomic versus atomless, BPI status relative to AC & E & must & Halmos 1963; Bell-Slomson 1969; Koppelberg 1989 \\
N-d-004 & P10 10.2 boolean-stone & Stone representation and duality: proof of the representation theorem, Stone spaces, clopen-algebra inverse, uniqueness of the countable atomless BA & T & must & Stone 1936; Halmos 1963; Johnstone 1982; Givant-Halmos 2009 \\
N-d-005 & P10 10.2 boolean-stone & Lindenbaum-Tarski algebras: algebraic completeness proof for CPL, free Boolean algebras & E & must & Tarski 1935; Rasiowa-Sikorski 1963 \\
N-d-006 & P10 10.2 boolean-stone & Priestley duality for distributive lattices & S & should & Priestley 1970; Davey-Priestley 2002 \\
N-d-007 & P10 10.2 boolean-stone & Stone-duality dividends for logic: compactness as compactness of the Stone space, type spaces of theories as Stone spaces & S & should & Johnstone 1982; Marker 2002 \\
N-d-008 & P10 10.3 heyting-algebras & Heyting algebras: definition, opens of a topological space, algebraic completeness of IPC, Rasiowa-Sikorski lemma & T & must & Heyting 1930; McKinsey-Tarski 1946; Rasiowa-Sikorski 1963; Balbes-Dwinger 1974 \\
N-d-009 & P10 10.3 heyting-algebras & Esakia duality and the lattice of intermediate logics survey & S & should & Esakia 1974; Chagrov-Zakharyaschev 1997 \\
N-d-010 & P10 10.3 heyting-algebras & locales and pointfree topology pointer & P & could & Johnstone 1982; Picado-Pultr 2012 \\
N-d-011 & P10 10.4 modal-algebras & Boolean algebras with operators: Jonsson-Tarski representation, canonical extensions, duality with general frames, algebraic completeness of normal modal logics & T & must & Jonsson-Tarski 1951; Goldblatt 1989; Blackburn-de Rijke-Venema 2001; Venema 2007 \\
N-d-012 & P10 10.4 modal-algebras & topological semantics: McKinsey-Tarski completeness of S4 for topological interior, completeness for the real line statement & T & should & McKinsey-Tarski 1944; Rasiowa-Sikorski 1963; van Benthem-Bezhanishvili 2007 \\
N-d-013 & P10 10.5 residuated-mv & MV-algebras and Chang completeness for infinite-valued Lukasiewicz logic; Mundici's l-group correspondence statement & T & must & Chang 1958; Chang 1959; Mundici 1986; Cignoli-D'Ottaviano-Mundici 2000 \\
N-d-014 & P10 10.5 residuated-mv & residuated lattices and FL-algebras: algebraization of substructural logics, subvariety lattice survey, cut elimination correspondences & S & must & Ono-Komori 1985; Galatos-Jipsen-Kowalski-Ono 2007 \\
N-d-015 & P10 10.5 residuated-mv & BL- and MTL-algebras and standard completeness statements & S & should & Hajek 1998; Cignoli-Esteva-Godo-Torrens 2000; Jenei-Montagna 2002 \\
N-d-016 & P10 10.6 algebraized-fol & cylindric algebras: algebraization of first-order logic, representability problems survey & S & should & Henkin-Monk-Tarski 1971; Henkin-Monk-Tarski 1985; Monk 2000 \\
N-d-017 & P10 10.6 algebraized-fol & relation algebras: Tarski's calculus of relations, representable relation algebras, undecidability statement, set theory without variables pointer & S & should & Tarski 1941; Tarski-Givant 1987; Maddux 2006 \\
N-d-018 & P10 10.6 algebraized-fol & Blok-Pigozzi algebraizability pointer & P & could & Blok-Pigozzi 1989; Font 2016 \\
N-d-019 & P10 10.7 category-toolkit & category theory for logicians: categories, functors, natural transformations, adjunctions, limits and colimits, Yoneda lemma statement & E & must & Mac Lane 1998; Awodey 2010 \\
N-d-020 & P10 10.7 category-toolkit & Lawvere algebraic theories and functorial semantics & S & should & Lawvere 1963; Adamek-Rosicky-Vitale 2011 \\
N-d-021 & P10 10.8 ccc-lambda & cartesian closed categories and typed lambda calculus: Curry-Howard-Lambek three-way correspondence with the equivalence theorem & T & must & Lambek 1980; Lambek-Scott 1986 \\
N-d-022 & P10 10.9 topos-overview & quantifiers as adjoints; hyperdoctrine view of first-order logic overview & S & should & Lawvere 1969; Pitts 2000 \\
N-d-023 & P10 10.9 topos-overview & elementary topos overview: subobject classifier, internal logic is intuitionistic HOL, presheaf examples, Kripke-Joyal semantics & S & must & Lawvere 1970; Lambek-Scott 1986; Mac Lane-Moerdijk 1992; Johnstone 2002 \\
N-d-024 & P10 10.9 topos-overview & topos-theoretic independence pointer: Tierney's sheaf-theoretic forcing and CH & P & could & Tierney 1972; Mac Lane-Moerdijk 1992 \\
N-x-001 & P7 -- & homotopy type theory and univalent foundations beyond a one-page pointer & P & -- & Univalent Foundations Program 2013 \\
N-x-002 & P7 -- & higher inductive types and cubical type theory internals & P & -- & Cohen-Coquand-Huber-Mortberg 2016 \\
N-x-003 & P6 -- & proper forcing, iterated forcing, and the forcing-axiom landscape (PFA/MM) beyond a pointer & P & -- & Shelah 1982; Foreman-Magidor-Shelah 1988 \\
N-x-004 & P6 -- & inner model theory, fine structure, and core models & P & -- & Jensen 1972; Mitchell-Steel 1994 \\
N-x-005 & P6 -- & PCF theory and advanced cardinal arithmetic beyond Silver's theorem statement & P & -- & Shelah 1994 \\
N-x-006 & P8 -- & quantum logic (orthomodular lattices as logic) & P & -- & Birkhoff-von Neumann 1936 \\
N-x-007 & P8 -- & adaptive logics & P & -- & Batens 2007 \\
N-x-008 & P11 -- & inductive logic and confirmation theory & P & -- & Carnap 1950 \\
N-x-009 & P8 -- & formal argumentation frameworks & P & -- & Dung 1995 \\
N-x-010 & P10 -- & abstract algebraic logic beyond the Blok-Pigozzi pointer & P & -- & Font-Jansana-Pigozzi 2003; Font 2016 \\
N-x-011 & P7 -- & geometry of interaction & P & -- & Girard 1989 \\
N-x-012 & P7 -- & proof nets beyond a pointer (correctness criteria, MLL net theory) & P & -- & Girard 1987; Danos-Regnier 1989 \\
N-x-013 & P10 -- & categorical logic beyond the CCC and topos overview (fibrations, sketches, infinity-topoi) & P & -- & Jacobs 1999; Lurie 2009 \\
N-x-014 & P4 -- & computable analysis and computability over the reals (TTE, Weihrauch degrees) & P & -- & Weihrauch 2000; Pour-El-Richards 1989 \\
N-x-015 & P9 -- & probabilistic and stochastic model checking, probability logics beyond a pointer & P & -- & Halpern 2003; Baier-Katoen 2008 \\
\end{longtable}
\endgroup
